% Additional extension; D actual-file and primary-source audit passed.

\begin{lemma}[Proved: from an SNC core to all exceptional valuations]
\label{D:Gorenstein-lc-valuative-tame}
Let $(S,s)$ be a Gorenstein log canonical surface singularity
over an algebraically closed field, and suppose that it is not
log terminal.  Every finite normal separable cover of $S$ that
is \'etale on the punctured germ is tame at every divisorial
valuation of $k(S)$ centred at $s$, including separability of
the corresponding residue extension.
\end{lemma}
\begin{proof}
By Lemma~\ref{D:Gorenstein-lc-configurations}, a minimal regular
resolution has either one smooth elliptic curve or a rational
nodal cycle as reduced exceptional fibre.  In the latter case
we blow up a possible self-node once.  We obtain a regular
resolution $Z\to S$ whose reduced exceptional divisor $E$ has
simple normal crossings.  Every cover in the statement is
tame along all components of $E$, by
Proposition~\ref{D:lc-elliptic-tame} or
Theorem~\ref{D:cusp-tame}.

We recall precisely the local persistence of this tameness.
Let $x\in E$, let $R=\mathcal O_{Z,x}^{\mathrm{sh}}$, and let
$f_1,\ldots,f_r$, with $r\leq2$, be the regular parameters
defining the boundary components through $x$.
The absolute Abhyankar lemma and its strictly local
corollary, \cite[XIII, Proposition 5.2 and Corollary 5.3,
pp.~428--430]{SGA1}, say that every connected component
of the punctured cover over $R$ is dominated by
\[
 R[T_1,\ldots,T_r]/
       (T_1^{n_1}-f_1,\ldots,T_r^{n_r}-f_r),
 \qquad n_i\text{ prime to the characteristic exponent}.
\]
The associated connected function-field extension is
Galois of prime-to-$p$ degree: all the required roots
of unity are present.  It is therefore tame at every
discrete valuation.  Every subextension is also tame.
This assertion descends from the strict localization:
an \'etale neighborhood of $x$ induces an unramified
extension at a valuation centred at $x$, and tameness
can be checked after such extensions.  One may equally
descend the finite Kummer domination to a finite
\'etale neighborhood and use the multiplicativity of
ramification indices and separability of residue
extensions.

For any specified divisorial valuation centred at $s$,
a further regular resolution extracts it.  A common
regular resolution with $Z$ is obtained from $Z$ by
point blowups, with centres in $E$ and its successive
inverse images.  The preceding local argument proves
tameness at each new exceptional divisor; the reduced
boundary remains simple normal crossings after every
such blowup.  It proves tameness at the specified
valuation, as required.  The tower and subextension
facts for tameness used here are, in particular,
\cite[Lemmas 15.116.5--9, section tag 0EXT]{StacksBIndexTame}.
\end{proof}

\begin{theorem}[Proved under stated hypotheses: non-klt points of prime-to-$p$ index]
\label{D:prime-to-p-index-tame}
Let $(S,s)$ be a normal log canonical surface singularity
over an algebraically closed field $k$, and suppose that it
is not log terminal.  Let $n$ be the local Cartier index of
$K_S$ at $s$.  If $n$ is prime to the characteristic exponent
of $k$, then every finite normal separable cover of $S$
that is \'etale on the punctured germ is tame at every
divisorial valuation centred at $s$.
\end{theorem}
\begin{proof}
We work in the strict henselization $A$ of the local ring
at $s$, as in Theorem~\ref{D:cusp-tame}; passing back to
the algebraic germ preserves the valuation conclusion.
The local index after this base change divides $n$, so
we may still use the given prime-to-$p$ integer $n$.
Write $U=\operatorname{Spec}A\setminus\{s\}$.
The canonical sheaf on $U$ is an invertible sheaf $L$,
and a local trivialization of $\omega_A^{[n]}$ gives
an isomorphism $L^{\otimes n}\simeq\mathcal O_U$.
Here square brackets denote the reflexive power.

\par\smallskip\noindent\textbf{Construction and properties of the cover.}\quad
The algebra
\[
 \mathcal O_U\oplus L^{-1}\oplus\cdots\oplus L^{-(n-1)},
\]
with multiplication specified by that trivialization,
defines a $\mu_n$-torsor over $U$.  Since $n$ is invertible,
it is finite \'etale.  Choose one connected component and
normalize $\operatorname{Spec}A$ in its function field
$K_1$.  Excellence makes the resulting map
$\pi:T\to\operatorname{Spec}A$ finite; it is normal and
quasi-\'etale, and its punctured open is the chosen torsor
component.  Because $k$ contains the $n$th roots of unity,
that component is a connected cyclic Galois cover of
degree $d\mid n$.  Explicitly, its stabilizer in the
constant group $\mu_n(k)$ is a cyclic group acting
simply transitively on each fibre of the component.
The covering normal ring is local, henselian, and has
residue field $k$.

On the torsor the line bundle $L$ has its tautological
trivialization.  Its pullback to every component is
therefore trivial.  Since the map is \'etale on $U$,
the canonical sheaf of the punctured $T$ is trivial.
Normality extends this equality of rank-one reflexive
sheaves across its omitted point, and gives
$\omega_T\simeq\mathcal O_T$.  A normal surface is
Cohen--Macaulay, so $T$ is Gorenstein.  In particular,
this conclusion does not require a hypothesis that
the reflexive cyclic algebra before normalization
was already regular or normal.

The cyclic extension $K_1/k(S)$ has prime-to-$p$ degree.
Its wild inertia is therefore trivial at every discrete
valuation.  Thus all its ramification indices are prime
to $p$, its residue extensions are separable, and the
different determinant has order $e-1$.  For every
divisorial valuation $F/E$ in this extension, the
canonical-divisor comparison yields
\begin{equation}\label{D:eq-index-discrepancy}
 a(F,T)+1=e(F/E)\bigl(a(E,S)+1\bigr).
\end{equation}
It follows that $T$ is log canonical.  Since $S$ is not
log terminal, some exceptional valuation $E$ has
discrepancy $-1$.  Every extension of it then has
discrepancy $-1$ by \eqref{D:eq-index-discrepancy}.
Thus $T$ is not log terminal either.  We have obtained
a Gorenstein non-klt log canonical cover without any
change of residue field.

\par\smallskip\noindent\textbf{Comparison for an arbitrary cover.}\quad
Let $M/k(S)$ be the finite separable extension of any
cover in the statement, and choose a compositum
$M_1=M K_1$.  Normalize $T$ in $M_1$.
It is a finite normal cover of $T$, \'etale on the
punctured germ, because this is true after base change
on that open.  By
Lemma~\ref{D:Gorenstein-lc-valuative-tame}, it is tame
at every exceptional divisorial valuation of $K_1$.
Every extension to $K_1$ of a divisorial valuation of
$k(S)$ centred at $s$ is still divisorial and is centred
at the closed point of $T$.

Consequently, at any such valuation the tower
\[
 k(S)\ \subset\ K_1\ \subset\ M_1
\]
is tame in both steps, and hence in its composite.
The subextension $M/k(S)$ is tame as well.
These implications include residue separability:
ramification indices multiply, and intermediate
fields of a finite separable residue extension are
separable.  See the opened proofs of
\cite[Lemmas 15.116.5 and 15.116.6]{StacksBIndexTame}.
This establishes the assertion for the original
arbitrary cover, not only for the cyclic index cover.
\end{proof}

\begin{corollary}[Proved: mixed omitted points with prime-to-$p$ lc index]
\label{D:prime-to-p-index-boundary}
Let $X$ be a normal proper integral surface over an
algebraically closed field, and let $Y\subset X$ be a
finite set of closed points.  Suppose that each point
of $Y$ is either log terminal, with no restriction on
its Cartier index, or log canonical with Cartier
index prime to the characteristic exponent.
Then $\pi_1(X\setminus Y)$ is topologically finitely
generated.  All singularities retained in the open
surface may be arbitrary.
\end{corollary}
\begin{proof}
Take a fixed regular proper resolution
$q:\widetilde X\to X$; its source is a projective
surface.  A finite \'etale cover of $X\setminus Y$
normalizes to a finite quasi-\'etale cover of $X$.
For every exceptional divisor above a non-klt point
of $Y$, Theorem~\ref{D:prime-to-p-index-tame} makes
the wild inertia image zero.  For each log terminal
point, the image is finite because it factors
through the finite punctured strict-local group
of \cite[Theorem 5.3, p.~739]{CarvajalRojasYasuda2025B}.
The fixed projective modification thus satisfies
Proposition~\ref{D:finite-exceptional-wild}.
As $Y$ has codimension two, that proposition gives
the asserted finite generation of the full group.
\end{proof}

\begin{remark}[Restriction on the index]
We have used the index hypothesis to obtain a separable
cyclic cover tame at every valuation and the exact
discrepancy equality \eqref{D:eq-index-discrepancy}.
When the index is divisible by the characteristic,
the cyclic construction can be inseparable, so the
\'etale canonical-pullback and tame-discrepancy
arguments just used do not apply.  This proof does
not establish preservation of log canonicity in
that case, and the preceding theorem makes no
assertion about it.  The failure of preservation
of \emph{log terminal} singularities by index-one
covers in small characteristic is explicitly
discussed in the introduction and Remark 5.5 of
\cite{CarvajalRojasYasuda2025B}.
\end{remark}
